\chapter{The BioDYM Workflow}
\label{ch:workflow}

The BioDYM workflow is organized into four sequential steps, each corresponding to a section in the Jupyter notebook.

\section{Overview}

\begin{enumerate}
  \item \textbf{Setup \& Data Loading} --- Load the Excel workbook, validate input data, initialize the MFA system
  \item \textbf{Calculation \& Validation} --- Run the iterative solver, check mass balance, verify results
  \item \textbf{Visualization} --- Explore results with interactive plots (Sankey, dynamics, stocks, composition)
  \item \textbf{Export \& Advanced Analysis} --- Export results, run Monte Carlo simulation, run scenarios
\end{enumerate}

\section{Step 1: Setup \& Data Loading}

The notebook begins by loading your Excel configuration file:

\begin{lstlisting}[style=python]
# Set path to your Excel input file
input_file = "01_data/01_input/your_system.xlsm"

# Load and validate
model_scope = define_model_scope(input_file)
mfa_system = initialize_mfa_system(model_scope)
mfa_system = define_flows_and_parameters(mfa_system, input_file)
\end{lstlisting}

During loading, BioDYM:
\begin{itemize}
  \item Reads all Excel sheets and validates their structure
  \item Creates Process, Flow, Stock, and Parameter objects
  \item Interpolates time series data to fill gaps between defined years
  \item Normalizes transfer coefficients so outflows sum to 100\% per process
  \item Calculates element compositions from the defined fractions
\end{itemize}

\begin{notebox}
If data points are only defined for certain years (e.g., 2020 and 2030), BioDYM automatically interpolates linearly for all intermediate years.
\end{notebox}

\section{Step 2: Calculation \& Validation}

The MFA calculation uses an iterative solver:

\begin{lstlisting}[style=python]
mfa_system_results = run_mfa_calculation(mfa_system, ...)
\end{lstlisting}

The solver repeatedly passes through all processes until no flow values change between iterations (convergence). In each pass:
\begin{enumerate}
  \item For each flow, the solver checks if the source process has calculated inflows
  \item If yes, it applies the appropriate process logic (Splitter, Transformer, DSM, FOMP)
  \item Transfer coefficients determine how inflow is distributed to outflows
  \item Final balances are calculated: $\Delta S = \text{inflow} - \text{outflow}$ for each process
\end{enumerate}

After calculation, the mass balance is displayed:

\begin{lstlisting}
Mass Balance Check (material):
  Process 1 (Forestry):        Error = 0.00e+00
  Process 2 (Processing):      Error = 0.00e+00
  Process 3 (Use Phase):       Error = 2.13e-10
\end{lstlisting}

Errors below $10^{-6}$ are normal (floating-point precision). Larger errors indicate a configuration problem.

\section{Step 3: Visualization}

BioDYM provides interactive Plotly-based visualizations (see Chapter~\ref{ch:visualization} for details):

\begin{itemize}
  \item \textbf{Sankey diagram} --- Overview of all flows, filterable by element and year
  \item \textbf{Process dynamics} --- Inflow, outflow, and stock change per process over time
  \item \textbf{Flow dynamics} --- Time series of individual flows
  \item \textbf{Stock analysis} --- Bar charts and time series of accumulated stocks
  \item \textbf{Composition analysis} --- Element composition of flows and stocks
  \item \textbf{DSM stock details} --- Per-category stock evolution for DSM processes
\end{itemize}

\section{Step 4: Export \& Advanced Analysis}

Results can be exported and further analyzed:

\begin{itemize}
  \item \textbf{Excel export} --- All flows, stocks, and parameters to \texttt{.xlsx}
  \item \textbf{KPI dashboard} --- Key performance indicators (total input, output, stock, recycling rates)
  \item \textbf{Monte Carlo} --- Uncertainty analysis (Chapter~\ref{ch:montecarlo})
  \item \textbf{Scenarios} --- Parameter comparison studies (Chapter~\ref{ch:scenarios})
\end{itemize}
